\section{The Square-Rule Passer Test: Definition, Checks and a Counterexample}\label{S-square}
This section defines the square-rule passer test S used for the repair in the main text, says what was computed with it, and shows that S is not a sound certificate in general. No proof is claimed.

\subsubsection*{1. Definitions}
\textbf{Board and game.} \texttt{k} files (cyclic on a cylinder, linear otherwise), \texttt{H} ranks. White pawns move up, towards rank \texttt{H}; Black pawns move down, towards rank 1. This is the frozen-king pawn game: a king move loses at once, a pawn that reaches its last rank wins at once, and a side with no legal pawn move loses.

\textbf{Passed (e.p.-aware).} A pawn of X on \texttt{(f, r)} is passed iff no enemy pawn lies on files \texttt{f-{}1, f, f+1} at a rank strictly beyond \texttt{r} in X's direction, and the pawn cannot be captured en passant.

\textbf{Runner.} A passed pawn is a runner iff no own pawn lies on the same file strictly beyond it.

\textbf{Distances.} For a White pawn on rank \texttt{r} let \texttt{sdist\_W(r) = H -{} r -{} (1 if r = 2 else 0)}, and for a Black pawn \texttt{sdist\_B(r) = r -{} 1 -{} (1 if r = H-{}1 else 0)}: the number of moves to promote if nothing is in the way, with one double step from the start rank. Let \texttt{n} be the least \texttt{sdist} over the runners of X, and \texttt{d} the least \texttt{sdist} over all pawns of the opponent Y (\texttt{d = \textrm{∞}} if Y has none).

\textbf{The test S.} S certifies X iff X has a runner and \texttt{n \textrm{≤} d} when X is to move, or \texttt{n \textless{} d} when Y is to move.

\subsubsection*{2. What was computed}
\begin{itemize}
\item \textbf{The repair of the main text.} In the sweep on the seven-file and six-file cylinders and the flat five-file board, all with 6 ranks, the game is ended only at a passed pawn that S certifies. This removes all 52 failures of the passer convention and leaves 128 mismatches, all from the kings (COMPUTED-1, \texttt{research/v2\_mechanisms/cc.py}). This is a comparison with the real results and does not depend on S being sound.
\item \textbf{Exhaustive check on small boards.} S never certifies a side that loses the frozen game in 158,544 cases (position, side to move, e.p. state): 1 and 2 files with 8 ranks and 2 files with 6 ranks, up to three pawns a side, and 3 files with 6 ranks, up to two pawns a side (COMPUTED-1, the same program).
\end{itemize}
\subsubsection*{3. S is not sound in general}
\textbf{Counterexample} (checked by hand and with the solver of \texttt{research/v2\_mechanisms/pawngame.py})\textbf{.} Flat board, 3 files and 8 ranks, \texttt{W: a4 b6 / B: c7}, Black to move. The pawn a4 is a runner with \texttt{n = 4}; Black's only pawn has \texttt{d = 5}; Black is to move and \texttt{4 \textless{} 5}, so S certifies White. But Black wins: 1... cxb6, 2. a5 (White's only pawn move), 2... bxa5, and White has no pawn move.

The reason is that an enemy pawn two files away from the runner can capture another pawn of X and land on a file next to the runner, ahead of it. The same position with 9 ranks (\texttt{W: a5 b7 / B: c8}) fails in the same way. The pattern needs a runner on rank \texttt{r}, a pawn of X on the next file at rank \texttt{r + 2} and an enemy pawn at rank \texttt{r + 3} on the file beyond, with \texttt{n \textless{} d} and Y to move. It needs three files and at least 8 ranks, and the exhaustive check above has at most 2 files at 8 ranks, so it does not meet it.

\textbf{Status.} S is a heuristic certificate. It is sound on the boards of the exhaustive check (COMPUTED-1). Whether it is sound on all boards with 6 ranks, and which extra condition makes it sound in general, is open. A sufficient extra condition must at least exclude enemy pawns that can reach a file next to the runner by capturing.

